%% Unit 1 -- An introduction to AI agents.
%% Ported from the standalone "AI Agents" lecture (KIT sdqbeamer deck, L2):
%% same layout and content, retyped against the tuftedark theme API.
%% No preamble here; \input by the master (disciplined_vibe.tex).
%%
%% Scientific sources are cited on the slide they support via \slidesources
%% (cite keys -> a numeric [n] marker up the right edge); the closing frame
%% prints this unit's cited papers with \printbibliography. Biblatex and the
%% per-unit refsection are set up in the master (disciplined_vibe.tex).

\providecommand{\afull}{$\bullet$}
\providecommand{\aemp}{\textcolor{tddim}{$\circ$}}

\sectionsubtitle{History, Theory and Mechanics}
% ======================================================================
\section[Foundations]{Foundations}
% ======================================================================

% ------------------------------------------------ software 1.0 / 2.0 / 3.0
\begin{frame}{The hottest new programming language is English}
  
  \vskip0.8em
  \begin{columns}[T,onlytextwidth]
    \begin{column}{0.31\textwidth}
      \textbf{Software 1.0}\par\smallskip
      {\small Humans write explicit instructions in code, and the machine
      executes them.}
    \end{column}
    \begin{column}{0.31\textwidth}
      \textbf{Software 2.0}\par\smallskip
      {\small Gradient descent writes the weights of neural networks from
      data.}
    \end{column}
    \begin{column}{0.31\textwidth}
      \textbf{Software 3.0}\par\smallskip
      {\small The program is a natural-language prompt, and the LLM
      compiles it to behaviour.}
    \end{column}
  \end{columns}
  \vskip1em
  \begin{tdblock}
    LLMs are a \alertbf{new computing substrate}, with prompts as programs and context windows as RAM.
  \end{tdblock}
\end{frame}


% ------------------------------------------------------- vibe coding: floor/bar
\begin{frame}{Raise the floor, keep the bar}
  \vskip0.8em
  {\small\textbf{Raise the floor.}\; Everyone can now build
  software they could not build before. Fast.}
  \begin{tdblock}
    State your intent, accept the diffs, run the result, embrace the exponentials and iterate
    without necessarily touching or reading every line. That working style is
    \alertbf{vibe coding}.
  \end{tdblock}
  \vspace{1.5em}
  {\small\textbf{Keep the bar.}\; AI engineering means you adjust your workbench to this abstraction layer.}
  \begin{tdblock}
    Most professional work requires rigor. As you will still own
    the code, the quality, and are accountable for whatever you ship.
    That working style is what we term \alertbf{discipline}.
  \end{tdblock}
\end{frame}




% ------------------------------------------------------------- quest: tokens
% \tok -- a single token chip, used to lay a word out as a sequence of tokens.
\newcommand{\tok}[1]{\tikz[baseline]{\node[draw=black!25,rounded corners,
  fill=white,inner sep=3pt,text=black,font=\ttfamily\small,anchor=base]{#1};}}
\begin{frame}{Making the machine understand words.}
  \slidesources{mikolov2013efficient,pennington2014glove,weizenbaum1966eliza,jurafsky2023slp}
  \begin{withmargin}
    \begin{tdmain}
      {\small\color{tddim}Word-level vocabulary}\\[0.4em]
      \tok{banana}\;$\rightarrow$\;\texttt{[24\,907]}\\[0.3em]
      {\footnotesize One token, one entry, one id. A larger vocabulary yields longer
      subwords, representing complex semantics and patterns directly.}
      \vskip1.2em
      {\small\color{tddim}Subword vocabulary}\\[0.4em]
      \tok{ba}\tok{na}\tok{na}\;$\rightarrow$\;\texttt{[3021, 512, 512]}\\[0.3em]
      {\footnotesize Three tokens from two entries, \texttt{ba} and \texttt{na}; the
      repeated \texttt{na} reuses id \texttt{512}. Fewer merges keep the vocabulary
      small and character-based, which improves generalization.}
    \end{tdmain}
    \begin{tdmargin}
      In the beginning was the \textcolor{DHBWred}{token}: a unit of text, a word
      or subword. It is the atomic unit a model reads and writes; a vocabulary is
      the fixed set of tokens it knows. Tokenizing is how text is sliced into
      entries of that set.
      \begin{tdblock}
        LLMs use vocabularies of 30{,}000+ tokens.
      \end{tdblock}
    \end{tdmargin}
  \end{withmargin}
\end{frame}


% ----------------------------------------------------------------- LLM basics
\begin{frame}{Making the machine understand language}
  \sideimage[0.42]{ghibli_parrot}
  \sidecaption{An LLM: the stochastic parrot that predicts the next token.}
  \slidesources{radford2019gpt2,devlin2018bert,raffel2019t5,touvron2023llama}
  \begin{sidebody}
    {\small
    \setlength{\abovedisplayskip}{4pt}\setlength{\belowdisplayskip}{4pt}%
    LLMs learn the structure and semantics of
    language by \alert{predicting the next token},

    \[\begin{aligned}
      \hat{y}_{t+1} &= \theta(x_{1:t}).\\
    \end{aligned}\]


    The free lunch, an implicit ground truth inside every
    token sequence, enables unsuperised learning of $\theta$ on unlabeled data, 

    \[\begin{aligned}
      \mathcal{D} &= \{(x_t, \tilde{y}_t)\}_{t=1}^{N},\\
      \textcolor{DHBWred}{x_{t+1}} &= \textcolor{DHBWred}{\tilde{y}_t}.
    \end{aligned}\]}
  \end{sidebody}
\end{frame}



% ------------------------------------------------------------------ attention
\begin{frame}{Attention is all you need}
  \slidesources{vaswani2017attention}
  \begin{withmargin}
    \begin{tdmain}
      Attention compares token vectors: the closer a query and a key point in the
      same direction, the more that token attends to the other, which reflects \alertbf{semantic similarity}.
      \vskip1.0em
      \centering
      \begin{tikzpicture}[>=Latex,scale=1.15]
        \coordinate (O) at (0,0);
        \coordinate (K) at (2.13,1.49);
        \coordinate (P) at (2.13,0);
        \draw[very thick,tdfg,->] (O) -- (3.2,0) node[right] {$\mathbf{Q}$};
        \draw[very thick,tdfg,->] (O) -- (K) node[above right] {$\mathbf{K}$};
        \draw[tddim] (0.72,0) arc (0:35:0.72);
        \draw[dashed,tddim] (K) -- (P);
        \draw[line width=1.6pt,DHBWred] (O) -- (P);
        \node[DHBWred,font=\footnotesize,below=2pt] at (1.06,0) {$\mathbf{Q}\mathbf{K}^\top$};
      \end{tikzpicture}
      \par\medskip
      {\small\color{tddim}The red projection is the dot product $\mathbf{Q}\mathbf{K}^\top$.}
    \end{tdmain}
    \begin{tdmargin}
      \textbf{Q}, query: what the token is looking for.
      \par\smallskip
      \textbf{K}, key: what a token offers to be matched against.
      \par\smallskip
      \textbf{V}, value: the information a token passes on once matched.
      \par\medskip
      Notice how these computations can be done in \alert{parallel}, 
      and accelerated by GPUs.
    \end{tdmargin}
  \end{withmargin}
\end{frame}



% ------------------------------------------------------------------ context
\begin{frame}{Context is all you have}
  \begin{withmargin}
    \begin{tdmain}
      Each token turns into a query and a key, and attention scores every one
      against every other, a grid of $\mathbf{Q}\mathbf{K}^\top$ interactions.
      This runs token by token over one flat sequence, the \alertbf{context}.
      \vskip1.1em
      \centering
      \begin{tikzpicture}[>=Latex,
          tokcell/.style={draw=tddim,rounded corners=2pt,minimum size=0.55cm,inner sep=0pt},
          faded/.style={draw=tddim!45}]
        \foreach \x in {0.62,1.24,1.86}{\node[tokcell,faded] at (\x,0){};}
        \foreach \x in {2.48,3.10,3.72,4.34,4.96,5.58,6.20}{\node[tokcell] at (\x,0){};}
        \draw[DHBWred,thick,rounded corners=3pt] (2.18,-0.35) rectangle (6.50,0.35);
        \node[DHBWred,font=\scriptsize] at (4.34,0.66) {context window};
        \node[tokcell,draw=DHBWred,fill=DHBWred,text=white,font=\scriptsize]
          (yn) at (7.35,0){$\hat{y}$};
        \draw[->,DHBWred,thick] (6.55,0) -- (yn.west);
      \end{tikzpicture}
      \par\medskip
      {\small\color{tddim}Context is limited - think Goldfish. Its size is capped by GPU memory.}
    \end{tdmain}
    \begin{tdmargin}
      Every token in the window competes for attention:
      \begin{itemize}
        \item your prompt,
        \item system prompts,
        \item the files you shared,
        \item tool and search results,
        \item the conversation so far,
        \item and the model's own replies.
      \end{itemize}
      \par\smallskip
      Irrelevant tokens still enter the $\mathbf{Q}\mathbf{K}^\top$ scores and
      crowd out the signal; doubling the tokens quadruples the compute.
    \end{tdmargin}
  \end{withmargin}
\end{frame}



% ------------------------------------------------------ prompting
\begin{frame}{Still a parrot: prompting}
  \slidesources{brown2020gpt3}
  \begin{withmargin}
    \begin{tdmain}
      \begin{block}{Prompt injection}
        Steer the model by engineering the input prompt, \alertbf{injecting instructions} or
        context that guide and nudge it toward a specific task without retraining.
      \end{block}
      \vskip0.6em
      This is zero-shot, the task lives entirely in the tokens you place in the
      context, so a good prompt is a good program.
    \end{tdmain}
    \begin{tdmargin}
      \textit{``The product is awesome is a good review.}\\
      \textit{The product breaks easily.''}\\
      $\rightarrow$ \textit{[LLM generates the response]}
    \end{tdmargin}
  \end{withmargin}
\end{frame}

% ------------------------------------------------------ instruction tuning
\begin{frame}{Still a parrot: Instruction tuning}
  \slidesources{wang2022selfinstruct,ouyang2022instructgpt,chung2022flan}
  \begin{withmargin}
    \begin{tdmain}
      \begin{block}{Instruction Behavior Training}
        \alertbf{Fine-tune} the model to follow an instruction. At this point the
        \textit{free lunch} is over. An \alert{instruct dataset} pairs each instruction with a label:
        Self-Instruct, InstructGPT, FLAN.
      \end{block}
      \vskip0.6em
      Prompting bends behaviour at inference; fine-tuning writes it into the weights.
      Prompt-to-tune is a recurrent pattern in LLM development.
    \end{tdmain}
    \begin{tdmargin}
      \textit{Instruction: ``Is this review positive?}\\
      \textit{The product breaks easily.''}\\
      $\rightarrow$ \textit{Response: ``No, it is negative.''}      
    \end{tdmargin}
  \end{withmargin}
\end{frame}







% ------------------------------------------------------------- quest: poem
\begin{frame}{\alertbf{Vibe} a poem on bananas}
  \sideimagecover[0.30]{banana_1}
  \sidecaption{The muse for your four-line banana poem.}
  \begin{sidebody}
    If not done yet, \alertbf{set up} your local LLM: clone the repo and run its bootstrap script.\par\smallskip
    {\small\href{https://github.com/Quillstacks/disciplinedvibe}{github.com/Quillstacks/disciplinedvibe}}
    \vspace{0.5em}
    \begin{itemize}
      \item Get comfortable with Open Code
      \item Write at least one poem, what works well, what does not?
      \item Share your best poem in the chat.
    \end{itemize}
    
    \vspace{1.5em}
    Not very useful outside the classroom, but it hints at the \alertbf{semantic capability} of LLMs and their speed, and
    at the applications they open up.   
  \end{sidebody}
  
\end{frame}






% References are collected in one shared package at the end of the deck
% (see disciplined_vibe.tex), not per unit.
